\section{Introduction}

The advent of sixth-generation (6G) wireless networks is anticipated to drive a paradigm shift from cell-centric architectures to user-centric Cell-Free Massive Multiple-Input Multiple-Output (CF-mMIMO) systems. By deploying a large number of distributed Access Points (APs) connected to a Central Processing Unit (CPU), CF-mMIMO eliminates cell-edge interference and provides uniform spectral efficiency across the coverage area \cite{Polese2023Understanding, Rajapaksha2023Unsupervised}. However, the practical realization of this architecture faces a critical impediment known as the fronthaul bottleneck. As the number of APs and antennas scales up to meet 6G density requirements, the volume of raw in-phase and quadrature (I/Q) samples required to be forwarded from APs to the CPU grows exponentially, overwhelming the limited capacity of standard fronthaul interfaces (e.g., eCPRI). Consequently, developing rigorous compression strategies that reduce fronthaul data rates without compromising signal detection accuracy has become a paramount engineering challenge.

Traditional approaches to fronthaul compression rely on uniform scalar quantization or vector quantization methods derived from rate-distortion theory, such as the Lloyd-Max algorithm. While these methods effectively minimize the Mean Squared Error (MSE) between the original and quantized signals, they are fundamentally agnostic to the downstream task. In a communication receiver, the ultimate goal is not to reconstruct the noisy received signal with high fidelity, but to correctly detect the transmitted user symbols. Allocating bits to reconstruct noise or interference in low Signal-to-Interference-plus-Noise Ratio (SINR) regimes is a waste of scarce fronthaul resources. Recently, the concept of Semantic Communication and Task-Oriented Communication has emerged as a promising alternative, advocating for the extraction and transmission of only the information relevant to the specific inference task \cite{Lan2021SemanticComm, Wen2023Integration}. This philosophy aligns with the Information Bottleneck (IB) principle, which seeks to maximize the mutual information between the compressed representation and the target variable (e.g., transmitted symbols) while minimizing the representation's complexity \cite{Shao2021IBApproach, Shao2022Multidevice}.

Leveraging deep learning, several works have explored task-oriented quantization for MIMO systems. Shlezinger \textit{et al.} proposed deep task-based quantization using hardware-limited scalar ADCs, demonstrating that learning the quantization mapping can approach optimal performance limits \cite{Shlezinger2021DeepTask}. Similarly, Zou \textit{et al.} introduced goal-oriented quantization for resource allocation, tailoring the quantization operation to the final optimization task \cite{Zou2022GoalOriented}. In the context of CF-mMIMO, recent studies have applied IB principles to relevance-based data compression, balancing informativity and compactness \cite{Danaee2024RelevanceBased, Danaee2025RelevanceBasedMultiUser}. Furthermore, Graph Neural Networks (GNNs) have been investigated for massive MIMO processing, with Feys \textit{et al.} utilizing GNNs and Gumbel-Softmax relaxation for energy-efficient downlink precoding \cite{Feys2025LearningToQuantize}. 

Despite these advancements, significant gaps remain. Most existing learning-based quantization schemes employ static bit-width allocations or decouple the quantization policy from the spatial aggregation mechanism. They often fail to dynamically adapt the quantization resolution of each AP on a per-symbol basis according to instantaneous channel conditions. Moreover, few works explicitly address the \textit{successive refinement} of signals in the uplink, where an AP should autonomously decide whether to transmit a coarse approximation, a refined residual, or nothing at all, based on the global interference landscape. The challenge lies in jointly optimizing the discrete, non-differentiable bit-width selection at distributed APs and the non-linear signal detection at the CPU.

To bridge this gap, this paper proposes a Task-Oriented Resource-Inference GNN (IB-CellFree-GNN) for uplink signal detection in backhaul-constrained CF-mMIMO systems. We treat quantization not as a signal reconstruction problem, but as a dynamic resource allocation problem. Our approach integrates a lightweight policy network at the edge with a powerful graph-based aggregator at the cloud, enabling end-to-end optimization of the rate-accuracy tradeoff. The main contributions of this paper are summarized as follows:

\begin{itemize}
    \item \textbf{Task-Oriented Split Architecture:} We propose a novel split-inference framework where distributed APs employ a lightweight Multi-Layer Perceptron (MLP) to compress local observations, while the CPU utilizes a GNN for global signal detection. This architecture shifts the computational burden of complex interference cancellation to the CPU while keeping edge processing minimal.
    
    \item \textbf{Differentiable Discrete Bit-Width Optimization:} We introduce a \textit{Resource-Inference Module} that utilizes Gumbel-Softmax relaxation to perform differentiable optimization of discrete bit-width choices (e.g., 0, 2, or 4 bits) during training. This allows the network to learn a policy that dynamically assigns higher precision to APs with favorable channel conditions and suppresses transmission from APs dominated by noise.
    
    \item \textbf{Bit-Aware Attention Mechanism:} We design a Bit-Aware GNN at the CPU that incorporates the quantization precision of each incoming message into the graph aggregation weights. This mechanism allows the CPU to robustly compensate for quantization noise by spatially filtering out low-precision contributions, effectively performing \textit{successive refinement} in the spatial domain.
    
    \item \textbf{Robustness and Efficiency:} Numerical results demonstrate that the proposed method achieves detection performance comparable to ideal full-precision (32-bit) systems while consuming an average of only 2.5 bits per AP. Furthermore, the system exhibits exceptional robustness, maintaining a bit error rate of approximately 2.1\% even under extreme conditions where 50\% of APs lose fronthaul connectivity, significantly outperforming fixed-bit baselines.
\end{itemize}

The remainder of this paper is organized as follows. Section II details the system model and the problem formulation. Section III presents the proposed IB-CellFree-GNN architecture, including the successive refinement quantization and the bit-aware graph aggregation. Section IV provides the numerical results and performance analysis. Finally, Section V concludes the paper.